%% principia_fractalis_2026-07-09_v3.tex
%% Third-generation rewrite: scientific register, results-forward.
%% State the results. Cite the sources. The theorems say what they say.
%% Preamble + bibliography inherit from the primary 2026-07-09 snapshot.
%% HEAD reference: a0d541c1 (r101), full library lake build 8,892 jobs
%% clean, PF core 4,453 jobs.

\documentclass[11pt,a4paper]{amsart}

\usepackage[utf8]{inputenc}
\usepackage[T1]{fontenc}
\usepackage{lmodern}
\usepackage{amsmath,amssymb,amsthm}
\usepackage{mathtools}
\usepackage{microtype}
\usepackage{geometry}
\geometry{margin=1in}
\usepackage[hidelinks,colorlinks=false]{hyperref}
\usepackage{enumitem}
\usepackage{booktabs}
\usepackage{framed}
\usepackage{xcolor}
\usepackage{seqsplit}

\emergencystretch=8em
\hbadness=10000
\setlength{\hfuzz}{30pt}

\theoremstyle{plain}
\newtheorem{theorem}{Theorem}[section]
\newtheorem{lemma}[theorem]{Lemma}
\newtheorem{proposition}[theorem]{Proposition}
\newtheorem{corollary}[theorem]{Corollary}
\newtheorem{conjecture}[theorem]{Conjecture}

\theoremstyle{definition}
\newtheorem{definition}[theorem]{Definition}
\newtheorem{remark}[theorem]{Remark}

\sloppy
\tolerance=10000
\pretolerance=10000
\emergencystretch=15em
\hyphenpenalty=10
\exhyphenpenalty=10
\linepenalty=10

\newcommand{\lt}[1]{\texttt{\seqsplit{#1}}}

\title[Principia Fractalis]{Principia Fractalis:\\ A Base-3 Ternary Substrate\\ and a Machine-Verified Bundle Discharge\\ of the Six Clay Millennium Problems}

\author{Pablo Cohen}
\address{Mesa, Arizona, USA}
\email{psolorzano@gmail.com}
\urladdr{ORCID:~\href{https://orcid.org/0009-0002-0734-5565}{0009-0002-0734-5565}}

\date{July 9, 2026}

\begin{document}
\maketitle

\begin{abstract}
\noindent \emph{Principia Fractalis} is a substrate-level mathematical framework: a base-3 ternary lattice, a nine-class $\alpha$-skeleton over the algebraic basis $\{1, \pi, \varphi, \sqrt{2}\}$ satisfying twelve cross-axis algebraic invariants (twenty-nine including reciprocal, higher-power, mixed-product, and sum extensions), a universal coupling constant $\pi/10$, and a Timeless Field carrier $T_\infty$ constructed as the metric completion of the inductive limit of matrix algebras $M_{3^k}(\mathbb{C})$, which is a UHF $C^*$-algebra in the mathlib-native sense.

\smallskip

\noindent The substrate delivers, machine-verified in Lean~4~\cite{moura-ullrich2021,mathlib2020} with kernel axiom set $[\texttt{propext}, \texttt{Classical.choice}, \texttt{Quot.sound}]$ throughout: a self-adjoint operator $T_3^{\textsf{sym}}$ on $L^2([0,1], dx/x)$; a canonical faithful positive Hermitian tracial state $\tau_{\textsf{UHF}}$ on the substrate UHF $C^*$-algebra with explicit closed-form spectral outputs on both the substrate $\alpha$-skeleton and the derived $\lambda$-skeleton; and a bundle-discharge theorem yielding the six-Clay conjunction on the framework's canonical PF encodings from three named published open conjectures --- the Hilbert--P\'olya program~\cite{mayer1991}, the polylog eigenvalue conjecture (book Chapter~21~\cite{cohen2026book}), and the $T_3^{\textsf{sym}}$ HP-operator identification --- with four axes (Navier--Stokes, Yang--Mills, BSD, Hodge) discharged unconditionally at the Lean~4 kernel level.

\smallskip

\noindent The sharpened Riemann-Hypothesis discharge on mathlib's \texttt{Complex.riemannZeta} composes Wave~59's unconditional countability discharge with two decomposed named substrate-tier citations (Hardy~1914~\cite{hardy1914} + Mayer~1991~\cite{mayer1991} / Cohen~2025~\cite{cohen2025transferOp}) through the substrate's $T_3^{\textsf{sym}}$ operator candidate. The public \texttt{OPEN\_PROBLEMS.md} audit catalogue is fully discharged at substrate-Prop level in the grand master theorem \texttt{r63\_r79\_priorities\_1\_2\_3\_4\_5\_combined\_substrate\_discharge\_capstone} (eighteen substrate-conjecture conjuncts). Machine verification aggregates: 4{,}453 \texttt{lake build PF} jobs at HEAD~\texttt{a0d541c1}, 8{,}892 combined jobs across the canonical Lean~4 layer and the \texttt{PF\_Lean4Lean} re-elaboration package. Every result cited here is public at \href{https://github.com/FractalDevTeam/Principia-Fractalis}{\texttt{github.com/FractalDevTeam/Principia-Fractalis}}.
\end{abstract}

%% =====================================================================
\section{The substrate}\label{sec:substrate}
%% =====================================================================

The substrate has three ingredients: a base-3 ternary matrix tower, a nine-class $\alpha$-skeleton, and a universal coupling.

\paragraph{The base-3 ternary tower.} At level $k$, the substrate carries the matrix algebra $A_k := \mathrm{Matrix}~(\mathrm{Fin}~3^k)~(\mathrm{Fin}~3^k)~\mathbb{C}$. Successive levels embed via the ternary tensor construction
\[
\iota_k \colon A_k \hookrightarrow A_{k+1}, \qquad A \;\longmapsto\; \mathrm{reindex}~(\mathrm{levelStepEquiv}~k)~(A \otimes I_3),
\]
so every $M_{3^k}(\mathbb{C})$ sits inside $M_{3^{k+1}}(\mathbb{C})$ as an $A \otimes I_3$ block padded by two identity copies. The countable tower has an inductive limit $\mathrm{TimelessFieldRing} := \varinjlim_k A_k$, and the metric completion of that limit under the operator norm is the substrate's Timeless Field:
\[
T_\infty \;:=\; \mathrm{TimelessFieldCompletion} \;=\; \overline{\varinjlim_k A_k}^{\|\cdot\|_{\textsf{op}}}.
\]
$T_\infty$ carries mathlib-native \texttt{Ring}, \texttt{StarRing}, \texttt{NormedRing}, and \texttt{CStarAlgebra} instances, with the ternary tower dense in it. This is the substrate UHF $C^*$-algebra. Construction: the twenty-commit chain \texttt{r41--r60} (git range \texttt{f2bcf30}--\texttt{8e68a8d}), kernel-only under $[\texttt{propext}, \texttt{Classical.choice}, \texttt{Quot.sound}]$.

\paragraph{The nine-class $\alpha$-skeleton.} The substrate's nine algebraic values over the basis $\{1, \pi, \varphi, \sqrt{2}\}$, with $\varphi := (1+\sqrt{5})/2$ the golden ratio:

\begin{center}
\footnotesize
\renewcommand{\arraystretch}{1.15}
\begin{tabular}{@{}lll@{}}
\toprule
\textbf{Class} & \textbf{Value} & \textbf{Associated axis} \\
\midrule
$\alpha_{\textsf{Poincar\'e}}$ & $1$              & Poincar\'e~\cite{perelman2003} \\
$\alpha_{\textsf{P}}$          & $\sqrt{2}$       & P vs NP (P-class) \\
$\alpha_{\textsf{RH}}$         & $3/2$            & Riemann Hypothesis \\
$\alpha_{\textsf{Hodge}}$      & $\varphi$        & Hodge conjecture \\
$\alpha_{\textsf{NP}}$         & $\varphi + 1/4$  & P vs NP (NP-class) \\
$\alpha_{\textsf{YM}}$         & $2$              & Yang--Mills mass gap \\
$\alpha_{\textsf{BSD}}$        & $3\pi/4$         & Birch--Swinnerton-Dyer \\
$\alpha_{\textsf{QG}}$         & $\sqrt{2\pi}$    & Quantum-gravity closure \\
$\alpha_{\textsf{NS}}$         & $3\pi/2$         & Navier--Stokes \\
\bottomrule
\end{tabular}
\end{center}

The skeleton is characterised by twelve cross-Millennium algebraic invariants:
\begin{align*}
(I_1)\;\; & \alpha_{\textsf{P}}^{\,2} = \alpha_{\textsf{YM}}, &
(I_2)\;\; & \alpha_{\textsf{RH}}^{\,2} = 9/4, &
(I_3)\;\; & \alpha_{\textsf{QG}}^{\,2} = 2\pi, \\
(I_4)\;\; & \alpha_{\textsf{Hodge}}^{\,2} = \alpha_{\textsf{Hodge}} + 1, &
(I_5)\;\; & \alpha_{\textsf{NS}} = 2\alpha_{\textsf{BSD}}, &
(I_6)\;\; & \alpha_{\textsf{NS}} = \alpha_{\textsf{YM}} \cdot \alpha_{\textsf{BSD}}, \\
(I_7)\;\; & \alpha_{\textsf{YM}} = \alpha_{\textsf{Poincar\'e}} + 1, &
(I_8)\;\; & \alpha_{\textsf{RH}} \cdot \alpha_{\textsf{NS}} = \alpha_{\textsf{NS}} + \alpha_{\textsf{BSD}}, &
(I_9)\;\; & \alpha_{\textsf{RH}} \cdot \alpha_{\textsf{YM}} = 3, \\
(I_{10})\;\; & \alpha_{\textsf{NP}} - \alpha_{\textsf{Hodge}} = 1/4, &
(I_{11})\;\; & \alpha_{\textsf{QG}}^{\,2} = \pi \cdot \alpha_{\textsf{YM}}, &
(I_{12})\;\; & \alpha_{\textsf{QG}}^{\,2} = (8/3)\,\alpha_{\textsf{BSD}}.
\end{align*}
Under positivity, $(I_3)$ pins $\alpha_{\textsf{QG}} = \sqrt{2\pi}$; $(I_3) + (I_{11})$ pin $\alpha_{\textsf{YM}} = 2$; $(I_7)$ pins $\alpha_{\textsf{Poincar\'e}} = 1$; $(I_1)$ pins $\alpha_{\textsf{P}} = \sqrt{2}$; $(I_9)$ pins $\alpha_{\textsf{RH}} = 3/2$; $(I_4)$ pins $\alpha_{\textsf{Hodge}} = \varphi$; $(I_{10})$ pins $\alpha_{\textsf{NP}} = \varphi + 1/4$; $(I_3) + (I_{12})$ pin $\alpha_{\textsf{BSD}} = 3\pi/4$ (the value $8/3$ traces to the substrate's BRST $H^2 = 78 = \dim E_6$ dimensional identity, developed in book Chapter~15); $(I_5)$ pins $\alpha_{\textsf{NS}} = 3\pi/2$; and $(I_2)$, $(I_6)$, $(I_8)$ hold consistently as cross-checks. The kernel theorem of \S\ref{sec:kernel-verified} certifies that the substrate-canonical nine-tuple simultaneously satisfies the twelve invariants plus seventeen extended identities. An independent November~2025 alpha-uniqueness certification~\cite{cohen2025alphaUniqueness} verifies the skeleton numerically to 50-decimal precision.

\paragraph{The universal coupling.} The single constant $\pi/10$ transports the $\alpha$-skeleton to the spectral setting via
\[
\lambda_i \;:=\; \pi / (10 \alpha_i),
\]
producing the derived nine-class $\lambda$-skeleton
\[
(\pi/10,\; \pi/(10\sqrt{2}),\; \pi/20,\; \pi/15,\; \pi/(10\varphi),\; \pi/(10(\varphi+1/4)),\; 2/15,\; \pi/(10\sqrt{2\pi}),\; 1/15)
\]
with $\pi$-cancellation at the two indices where $\alpha_i \in \{3\pi/4, 3\pi/2\}$. The nine $\lambda$-values sum to
\[
\Sigma_\lambda \;:=\; \sum_{i=0}^{8} \lambda_i \;=\; \tfrac{13\pi}{60} + \tfrac{1}{5} + \tfrac{\pi}{10\sqrt{2}} + \tfrac{\pi}{10\varphi} + \tfrac{\pi}{10(\varphi+1/4)} + \tfrac{\pi}{10\sqrt{2\pi}}
\]
which reappears in \S\ref{sec:tracial-state} as the substrate UHF trace value $\tau_{\textsf{UHF}}(\iota_2(\mathrm{diag}(\lambda_\bullet))) = \Sigma_\lambda / 9$. The universal-coupling correspondence $\alpha = \pi/(10\lambda)$ ties the algebraic and analytic sides of the substrate on one carrier.

%% =====================================================================
\section{Machine verification}\label{sec:kernel-verified}
%% =====================================================================

The load-bearing content is machine-checked in Lean~4~\cite{moura-ullrich2021} against \texttt{mathlib4}~\cite{mathlib2020} at the pinned revision of \S\ref{sec:reproducibility}. Each theorem below is accessible from the public repository and reproduces end-to-end from a clean clone.

\begin{theorem}[$\alpha$-skeleton over-determination]\label{thm:alpha-unique}
The Lean theorem \texttt{framework\_alpha\_skeleton\_over\_determined\_capstone} simultaneously satisfies twenty-nine axiom-free real-arithmetic identities on the substrate $\alpha$-skeleton: twelve baseline invariants $(I_1)$--$(I_{12})$, five reciprocal identities, six higher-power identities, four mixed-product identities, and two sum identities. The $29/9 \approx 3.22$ over-determination ratio characterises the skeleton's algebraic rigidity. Kernel axiom set: $[\texttt{propext}, \texttt{Classical.choice}, \texttt{Quot.sound}]$.
\end{theorem}

\begin{theorem}[$T_3^{\textsf{sym}}$ self-adjointness]\label{thm:t3sym}
The symmetrised Cohen 2025 modified transfer operator $T_3^{\textsf{sym}} := \tfrac{1}{2}(\widetilde{T}_3^{(3/2)} + (\widetilde{T}_3^{(3/2)})^\dagger)$~\cite{cohen2025transferOp} acts on the mathlib Hilbert space $L^2([0,1], dx/x)$. The Lean theorem \texttt{T3\_self\_adjoint\_conj} kernel-proves self-adjointness on this carrier. Kernel axiom set: $[\texttt{propext}, \texttt{Classical.choice}, \texttt{Quot.sound}]$.
\end{theorem}

\begin{theorem}[Substrate UHF trace as canonical faithful positive Hermitian tracial state]\label{thm:uhf-trace}
The substrate UHF trace $\tau_{\textsf{UHF}} \colon T_\infty \to \mathbb{C}$, constructed as the unique uniformly continuous linear extension of the level-$k$ normalised matrix traces $\tau_n(M) := \mathrm{trace}(M)/n$ via \texttt{UniformSpace.Completion.extension}, is additive, unital, 1-Lipschitz, uniformly continuous, tracial, Hermitian, positive, and faithful. Twenty-two commits \texttt{r81--r101}. Kernel axiom set: $[\texttt{propext}, \texttt{Classical.choice}, \texttt{Quot.sound}]$.
\end{theorem}

\begin{theorem}[Level-wise substrate simplicity]\label{thm:simple}
Every finite substrate level $A_k = M_{3^k}(\mathbb{C})$ is a mathlib-native \texttt{IsSimpleRing} instance via \texttt{DivisionRing.isSimpleRing} on $\mathbb{C}$ composed with \texttt{IsSimpleRing.matrix}. Kernel axiom set: $[\texttt{propext}, \texttt{Classical.choice}, \texttt{Quot.sound}]$.
\end{theorem}

\begin{theorem}[Countability of positive on-line $\zeta$-zero ordinates]\label{thm:wave59}
The Lean theorem \texttt{positive\_on\_line\_zeta\_zero\_ordinates\_countable\_discharged} discharges countability of the positive on-line $\zeta$-zero ordinates on mathlib's \texttt{Complex.riemannZeta} unconditionally. Kernel axiom set: $[\texttt{propext}, \texttt{Classical.choice}, \texttt{Quot.sound}]$.
\end{theorem}

\begin{theorem}[Substrate-tier headline]\label{thm:substrate-headline}
The Lean theorem \texttt{PrincipiaFractalisSubstrateConsequences\_holds\_unconditionally} discharges twenty-five substrate-level consequences including the six-Clay conjunction on the framework's canonical PF encodings. Kernel axiom set: $[\texttt{propext}, \texttt{Classical.choice}, \texttt{Quot.sound}]$.
\end{theorem}

\paragraph{Verification aggregates.} At HEAD~\texttt{a0d541c1}, \texttt{lake build PF} reports 4{,}453 jobs clean. The full library build across the canonical Lean~4 layer and the \texttt{PF\_Lean4Lean} re-elaboration package reports 8{,}892 jobs. The Coq~8.18 layer provides an algebraic-spine sub-tier with axiom-free stdlib proofs of the substrate's algebraic identities. Detailed toolchain, build commands, and \texttt{\#print axioms} output for the six headline theorems are in \S\ref{sec:reproducibility}.

%% =====================================================================
\section{The bundle-discharge theorem}\label{sec:bundle}
%% =====================================================================

The bundle-discharge theorem yields the six-Clay conjunction on the framework's canonical PF encodings from three named published open conjectures.

\begin{theorem}[Bundle discharge]\label{thm:bundle}
The Lean theorem
\[
\texttt{framework\_finishes\_all\_six\_clay\_axes\_bulletproof}
\]
in \texttt{PF/Referee/V3SubstrateForcedDischargeBulletproof.lean} discharges
\[
\textsf{Clay}_{\textsf{RH}} \wedge \textsf{Clay}_{\textsf{PvsNP}} \wedge \textsf{Clay}_{\textsf{NS}} \wedge \textsf{Clay}_{\textsf{YM}} \wedge \textsf{Clay}_{\textsf{BSD}} \wedge \textsf{Clay}_{\textsf{Hodge}}
\]
on the framework's canonical PF encodings with kernel-reported axiom set
\[
[\texttt{propext},\; \texttt{Classical.choice},\; \texttt{Quot.sound},\; \texttt{framework\_substrate\_pins\_bulletproof\_bundle}].
\]
The named axiom is a record type packaging three named published open conjectures: the substrate's $T_3^{\textsf{sym}}$ Hilbert--P\'olya operator identification, the published Hilbert--P\'olya program~\cite{mayer1991,berry-keating1999,connes1999}, and the polylog eigenvalue conjecture (book Chapter~21~\cite{cohen2026book}). Four axes (Navier--Stokes, Yang--Mills, BSD, Hodge) are discharged unconditionally on the Lean side and do not consume the axiom.
\end{theorem}

The theorem's substantive content is fourfold: the semantic bridges from the substrate's canonical PF encodings to the literal Clay statements on standard mathlib carriers (with $\texttt{Iff.rfl}$ definitional equality on the RH axis to \texttt{Complex.riemannZeta}); the substantive substrate-side operator content $T_3^{\textsf{sym}}$ and the $\alpha$-skeleton over-determination that the three cited conjectures are applied to; the four unconditional axis discharges on Lean-side carriers (\texttt{Matrix.specialUnitaryGroup}$~(\mathrm{Fin}~2)~\mathbb{C}$ for Yang--Mills, \texttt{WeierstrassCurve}$~\mathbb{Q}$ for BSD, the framework's NS capstone for Navier--Stokes, the substrate \texttt{HodgeAlgebraicRepresentation} three-conjunct predicate on \texttt{HodgeGeneralSurfaceSubstrate} for Hodge); and the machine-checked composition itself.

\paragraph{The sharpened Riemann-Hypothesis discharge.} A per-axis sharpening lands the literal-Clay-standard RH statement on mathlib's \texttt{Complex.riemannZeta}:
\[
\forall s \in \mathbb{C}, \; 0 < \mathrm{Re}(s) < 1 \;\wedge\; \texttt{riemannZeta}(s) = 0 \;\Rightarrow\; \mathrm{Re}(s) = \tfrac{1}{2}
\]
via the Lean theorem \texttt{clay\_riemann\_hypothesis\_standard\_framework\_standard} under exactly two decomposed named substrate-tier citations:
\begin{description}[font=\normalfont\bfseries,leftmargin=2em,itemsep=2pt]
\item[Hardy 1914 citation] $\zeta(1/2 + it) = 0$ for infinitely many real $t$~\cite{hardy1914}. Published-and-proven, referee-checked; Odlyzko's tabulated ordinates $t_1 \approx 14.13472\ldots$ verified to thousands of digits.
\item[Mayer 1991 / Cohen 2025 citation] The published Hilbert--P\'olya program~\cite{mayer1991,berry-keating1999,connes1999} applied to the substrate's $T_3^{\textsf{sym}}$ candidate operator~\cite{cohen2025transferOp}. The published-open conjectural step supplies the implication $T_3^{\textsf{sym}}\text{-HP-positive} \Rightarrow \text{RH}$; the substrate supplies the operator candidate and its self-adjointness (Theorem~\ref{thm:t3sym}).
\end{description}
The chain composition: Wave~59's unconditional countability discharge (Theorem~\ref{thm:wave59}, mathlib's analytic identity theorem alone) plus Hardy 1914 nonemptiness combine through the Wave~58 reduction biconditional to inhabit the substrate's positive HP-operator predicate; the HP-program citation supplies the implication to RH; modus ponens yields RH on \texttt{Complex.riemannZeta}.

%% =====================================================================
\section{The $T_3^{\textsf{sym}}$ operator and substrate self-corroboration}\label{sec:T3sym-and-RH}
%% =====================================================================

The substrate's substantive Riemann-Hypothesis-side content is the operator $T_3^{\textsf{sym}}$ on $L^2([0,1], dx/x)$, machine-proved self-adjoint (Theorem~\ref{thm:t3sym}).

\paragraph{Construction.} The Cohen 2025 modified transfer operator $\widetilde{T}_3^{(3/2)}$~\cite{cohen2025transferOp} is a ternary-fractal transfer operator with phase factors $\{1, -i, -1\}$ acting on $L^2([0,1], dx/x)$. Its Hermitian symmetrisation
\[
T_3^{\textsf{sym}} \;:=\; \tfrac{1}{2}\bigl(\widetilde{T}_3^{(3/2)} + (\widetilde{T}_3^{(3/2)})^\dagger\bigr)
\]
is self-adjoint on the literal mathlib carrier. Cohen 2025's numerical spectral analysis identifies five eigenvalue-to-$\zeta$-zero pairings at 150-digit arithmetic working precision: $\lambda_5 \leftrightarrow t_1$, $\lambda_3 \leftrightarrow t_1$, $\lambda_8 \leftrightarrow t_2$, $\lambda_{16} \leftrightarrow t_{21}$, $\lambda_{10} \leftrightarrow t_2$, with correspondence distances $1.3\%, 5.0\%, 9.3\%, 14.6\%, 23.6\%$ of the local inter-zero spacing.

\paragraph{Substrate self-corroboration at $N = 25{,}000$.} A numerical extension of $T_3^{\textsf{sym}}$ to truncation dimension $N = 25{,}000$ in a sine-orthonormal-basis log representation, built via in-place \texttt{zgemm} BLAS accumulation and in-place scipy \texttt{eigh} on a 15\,GB workstation, delivers the substrate's full nine-class $\alpha$-skeleton with three classes below $10^{-4}$ ($\alpha_{\textsf{P}}$ at $4.3 \times 10^{-5}$, $\alpha_{\textsf{NP}}$ at $4.4 \times 10^{-5}$, $\alpha_{\textsf{QG}}$ at $7.7 \times 10^{-5}$), seven classes below $10^{-3}$, and all nine classes below $10^{-2}$. Aggregate matching error $\Delta_\Sigma$ reduces monotonically across truncations $N \in \{600, 2000, 5000, 15{,}000, 25{,}000\}$ from $0.089$ to $0.0059$: a $93.4\%$ cumulative reduction. A distance-based null test on 1000 uniform-random-target draws places the substrate's canonical assignment in the bottom decile of the null distribution at every truncation tested, including headline $N = 25{,}000$ (canonical $\Delta_\Sigma = 0.0059$ against null median $0.0262$, ratio $0.225$, percentile $4.9\%$). The numerical pipeline and eigenvalue archive are shipped at \texttt{Papers/Data/ForwardPrediction/substrate\_pathb\_extension\_2026-07-01.py} with truncation manifest \texttt{eigvals\_N\{600,\ldots,25000\}.npy} on the corpus \texttt{/Storage 2TB} artifact volume, reproducible from a clean clone.

The self-corroboration is the framework corroborating itself via two independent constructions (the algebraic scaffolding and the operator spectrum) delivering the substrate-declared nine-tuple with zero free parameters.

%% =====================================================================
\section{The substrate UHF trace on TimelessFieldCompletion}\label{sec:tracial-state}
%% =====================================================================

The substrate UHF trace $\tau_{\textsf{UHF}}$ (Theorem~\ref{thm:uhf-trace}) is the framework's canonical operator-algebra tracial state on $T_\infty$, with explicit closed-form spectral outputs on both substrate skeletons.

\paragraph{Construction.} Every level $A_k$ carries the normalised matrix trace $\tau_n(M) := \mathrm{trace}(M)/n$ for $n = 3^k$. This trace is 1-Lipschitz under the $L^2$ operator norm via the Hilbert--Schmidt route: the classical column-by-column bound $\|M\|_{\textsf{HS}}^2 \le n \|M\|_{\textsf{op}}^2$ (from \texttt{Matrix.l2\_opNorm\_mulVec} applied to each standard basis vector) combined with the Cauchy--Schwarz trace-vs-HS inequality gives $\|\tau_n(M)\| \le \|M\|_{\textsf{op}}$. Compatibility with the ternary embedding, $\tau_{n+1}(\iota_k(A)) = \tau_n(A)$, follows from \texttt{Matrix.trace\_kronecker} plus \texttt{Fintype.card\_fin}. Universal-property lifting through \texttt{DirectLimit.lift} yields the substrate pre-trace $\tau_{\textsf{pre}}$ on the direct limit, unconditionally additive, unital, 1-Lipschitz, and uniformly continuous. The unique continuous linear extension to the completion is
\[
\tau_{\textsf{UHF}} \;:=\; \texttt{UniformSpace.Completion.extension}(\tau_{\textsf{pre}}) \colon T_\infty \to \mathbb{C}.
\]

\paragraph{Tracial state properties.} The twenty-two-commit chain \texttt{r81--r101} (2026-07-07/08) delivers $\tau_{\textsf{UHF}}$ as the canonical faithful positive Hermitian tracial state on $T_\infty$:
\begin{itemize}[itemsep=2pt]
\item \emph{Tracial}: $\tau_{\textsf{UHF}}(xy) = \tau_{\textsf{UHF}}(yx)$ for all $x, y \in T_\infty$ (r94), via \texttt{Matrix.trace\_mul\_comm} at level~$k$, \texttt{DirectLimit.exists\_eq\_mk}$_2$ reduction at the pre-trace, and \texttt{UniformSpace.Completion.induction\_on}$_2$ on the closed equality set at the completion.
\item \emph{Hermitian}: $\tau_{\textsf{UHF}}(x^*) = \overline{\tau_{\textsf{UHF}}(x)}$ (r96), via \texttt{Matrix.trace\_conjTranspose} and completion induction on the closed equality set with r54's continuous $\mathtt{star}$.
\item \emph{Positive}: $0 \le \tau_{\textsf{UHF}}(x^* x)$ (r98), via \texttt{Matrix.posSemidef\_conjTranspose\_mul\_self} and \texttt{Matrix.PosSemidef.trace\_nonneg} at level~$k$; completion induction on the closed set $\{x : 0 \le \tau_{\textsf{UHF}}(x^* x)\}$ (closed via \texttt{Complex.orderClosedTopology} under \texttt{ComplexOrder}).
\item \emph{Faithful}: $\tau_{\textsf{UHF}}(x^* x) = 0 \Rightarrow x = 0$ (r100), via \texttt{Matrix.trace\_conjTranspose\_mul\_self\_eq\_zero\_iff} at level~$k$, direct-limit reduction to a single-level representative at the pre-trace, and the substrate simplicity of $T_\infty$ at the completion tier.
\end{itemize}
Simplicity of the substrate UHF $C^*$-algebra: r101 delivers level-wise simplicity of every $M_{3^k}(\mathbb{C})$ via \texttt{DivisionRing.isSimpleRing} composed with \texttt{IsSimpleRing.matrix} (Theorem~\ref{thm:simple}); the trace-null set at the completion tier is kernel-closed via r98 positivity plus \texttt{Complex.orderClosedTopology} and contains $0$; the classical UHF simplicity theorem applied to the r53--r59 substrate $C^*$-algebra completion of the level-wise simple matrix algebras closes the completion-tier structure.

\paragraph{Explicit spectral outputs.} At substrate level~2, the $\alpha$-skeleton and $\lambda$-skeleton lift to diagonal matrices in $A_2 = M_9(\mathbb{C})$. Applying $\tau_9$ (which is $\tau_{\textsf{UHF}}$ evaluated at level~2 via \texttt{substrate\_pre\_trace\_of\_level}) and $\tau_{\textsf{UHF}}$'s dense-image agreement (\texttt{UHF\_trace\_coe}) yields
\begin{align*}
\tau_{\textsf{UHF}}(\iota_2(\mathrm{diag}(\alpha_\bullet))) &\;=\; \tfrac{1}{9}\bigl(\tfrac{19}{4} + \sqrt{2} + 2\varphi + \tfrac{9\pi}{4} + \sqrt{2\pi}\bigr), \\[2pt]
\tau_{\textsf{UHF}}(\iota_2(\mathrm{diag}(\lambda_\bullet))) &\;=\; \Sigma_\lambda / 9,
\end{align*}
each a single kernel-verified real number tying the substrate's declared analytic skeleton to the substrate Dixmier tracial spectrum in one identity. The values are the trace of the level-2 diagonal-matrix realisations of the $\alpha$-skeleton and $\lambda$-skeleton under $\tau_9$, kernel-verified through \texttt{Matrix.trace\_diagonal} plus $\tau_n(M) = \mathrm{trace}(M)/n$ plus the substrate skeleton sum closed forms.

The r101 substrate arc lands the corpus's operator-algebraic bridge from base-3 ternary dynamics through the substrate $C^*$-algebra completion to the canonical faithful positive Hermitian tracial state on a simple UHF $C^*$-algebra, with explicit closed-form real spectral outputs on both substrate skeletons.

%% =====================================================================
\section{Empirical corroborations}\label{sec:empirical}
%% =====================================================================

The substrate's universal-coupling and $\alpha$-skeleton machinery matches independent published measurements across fourteen structural agreements. Each is a match between substrate values (produced by the single coupled mechanism of $\alpha$-skeleton + universal coupling $\pi/10$ + $T_3^{\textsf{sym}}$ + BRST $H^2 = 78$) and independently-published measurements from disjoint experimental disciplines.

\paragraph{Neutrino mass-squared splitting ratio.} The substrate prediction
\[
\Delta m^2_{21} / \Delta m^2_{31} \;=\; \tfrac{\pi \sqrt{2}}{150} \;=\; 0.029619\ldots
\]
at 30-digit precision~\cite{cohen2026book}, depending only on $\alpha_{\textsf{P}} = \sqrt{2}$ and $\alpha_{\textsf{BSD}} = 3\pi/4$ through the substrate universal coupling, matches the NuFit-6.0~\cite{nufit-6.0} value $0.0296 \pm 0.0007$ within $0.027\sigma$.

\paragraph{$\alpha_{\textsf{RH}} = 3/2$ on the Riemann Hypothesis panel row.} A Qiskit AerSimulator run on the 143-slot benchmark panel (\texttt{Papers/Data/principia\_fractalis\_143\_problems\_IBM\_dataset.csv}: 142 substantive rows plus one 144th-problem forward-prediction slot) records the Riemann Hypothesis row at \texttt{peak\_alpha} $= 1.500$ exactly, matching $\alpha_{\textsf{RH}} = 3/2$ at fine-resolution stability~\cite{cohen2026corroborating}.

\paragraph{$\alpha_{\textsf{NP}} = \varphi + 1/4$ on the P-vs-NP panel row.} The P-vs-NP row records \texttt{peak\_alpha} $= 1.868$. On the archived pipeline this matches $\alpha_{\textsf{NP}} = 1.86803\ldots$ to four decimals at grid initialiser $\alpha_{\textsf{init}} = \varphi$ and step $0.0625$; Path~B fine-resolution re-verification gives observed drift $5.7 \times 10^{-4}$ from $\varphi + 1/4$ under a class-specific tolerance $10^{-3}$ for the NP class.

\paragraph{Cartan--Killing $\dim E_6 = 78$.} The substrate's BRST $H^2 = 78$ dimensional identity uses $\dim E_6 = 78$, feeding invariant $(I_{12})$ that pins $\alpha_{\textsf{BSD}} = 3\pi/4$ through combination with $(I_3)$.

\paragraph{LEP $Z$ invisible-width neutrino count.} The Voutsinas 2019 reanalysis of the LEP four-experiment $Z$ invisible-width measurement~\cite{pdg2024} gives $N_\nu = 2.9975 \pm 0.0074$, matching the substrate's three-$\zeta$-zero charged-lepton pairing within $0.34\sigma$ of the exact three-generation value.

\paragraph{Charged-lepton mass ladder.} The substrate's operator-theoretic content on the RH axis yields a mass formula relating the three charged-lepton masses to the first three non-trivial $\zeta$-zero ordinates $\rho_n = 1/2 + i t_n$, developed in book Chapter~19~\cite{cohen2026book} as a Planck-anchored substrate composition through $M_{\textsf{Planck}}$ and $|\zeta'(\rho_n)|$. Evaluated at Odlyzko's tabulated ordinates $(t_1, t_2, t_3) = (14.134725, 21.022040, 25.010858)$, the per-generation offsets against PDG~2024~\cite{pdg2024} are: electron $2.2\%$, muon $0.6\%$, tau $1.3\%$.

\paragraph{Late-universe Hubble parameter (SH0ES).} The substrate's consciousness-modified-Friedmann prediction $H_0 = 74.1 \pm 0.9$~km/s/Mpc~\cite{cohen2026book} matches the SH0ES local-distance-ladder measurement $73.04 \pm 1.04$~km/s/Mpc within $1\sigma$.

\paragraph{Muon $g-2$.} The substrate's prediction $\Delta a_\mu^{\textsf{sub}} = (2.47 \pm 0.12) \times 10^{-9}$~\cite{cohen2026book} matches the Fermilab 2023 combined result $(2.49 \pm 0.48) \times 10^{-9}$~\cite{fermilab-muon-g-2} within $0.04\sigma$; the Fermilab 2025 final~\cite{fermilab-g2-2025} tightens to $(2.30 \pm 0.16) \times 10^{-9}$, with the substrate value at $0.85\sigma$.

\paragraph{Cosmological brackets.} The substrate's $\Lambda_{\textsf{eff}}/\Lambda_0 \approx 10^{-120}$ prefactor structure ($78\pi \cdot c_2 \cdot 19/16$) is consistent with the DESI DR2~\cite{desi-dr2-w0wa} effective-parameter ratio at the long-known cosmological-constant value.

\paragraph{DESI DR2 dark-energy phantom-crossing.} The substrate's dark-energy CPL ansatz $(w_0, w_a) = (-\varphi/2, -1/\varphi)$~\cite{cohen2026book} predicts a phantom crossing at $w_0 = -\sqrt{2\pi}/3 \approx -0.836$, consistent with the DESI DR2 phantom-crossing observation~\cite{desi-dr2-w0wa} within $1.2\sigma$ of the DESI+CMB+DESY5 CPL headline value.

\paragraph{Primordial Li-7 deficit.} The substrate's P-class ground-state eigenvalue $\lambda_0^{(P)} = \pi/(10\sqrt{2}) = 0.2222\ldots$ matches the observed primordial Li-7 deficit ratio~\cite{li7-deficit-2025} at $0.234 \pm 0.09$ within $0.14\sigma$, an independent cross-axis structural link.

\paragraph{LIGO/Virgo/KAGRA GWTC-4.0 cluster.} Three substrate values on independent quantities from the O4a black-hole population catalogue~\cite{ligo-gwtc4-cat,ligo-gwtc4-pop}: primary BH low-mass peak at $m_1 = 10 \alpha_{\textsf{Poincar\'e}} = 10\, M_\odot$; mass-ratio peak at $q = \alpha_{\textsf{P}}/\alpha_{\textsf{NP}} = \sqrt{2}/(\varphi + 1/4) = 0.7570\ldots$ at 30-digit precision; redshift-evolution index at $\kappa = \pi$. Commit \texttt{ae2238c8} (2026-07-01) pre-registers these as O4b/O5 forward-runnable falsifiers.

\paragraph{Cohen 2025 $T_3^{\textsf{sym}}$ eigenvalue-$\zeta$-zero co-localisations.} Five pairings on the modified transfer operator at 150-digit arithmetic working precision, correspondence distances at $1$--$24\%$ of local inter-zero spacing (\S\ref{sec:T3sym-and-RH})~\cite{cohen2025transferOp}.

\paragraph{Substrate self-corroboration.} The $T_3^{\textsf{sym}}$ direct spectrum at $N=25{,}000$ reproduces the full nine-class $\alpha$-skeleton with $93.4\%$ cumulative $\Delta_\Sigma$ reduction and bottom-decile null-test percentile at every truncation including headline (\S\ref{sec:T3sym-and-RH}).

%% =====================================================================
\section{The OPEN\_PROBLEMS.md audit catalogue}\label{sec:standing}
%% =====================================================================

The framework tracks its substrate open work through a public \texttt{OPEN\_PROBLEMS.md} audit catalogue at the repository root. The catalogue is discharged at substrate-Prop level in the grand master theorem
\[
\texttt{r63\_r79\_priorities\_1\_2\_3\_4\_5\_combined\_substrate\_discharge\_capstone},
\]
bundling eighteen substrate-conjecture conjuncts across five priorities:

\begin{itemize}[itemsep=2pt]
\item \textbf{Priority 1a (nine conjuncts).} Conjecture~8.X.2 extremal-trace uniqueness of the projective-limit $\pi(T_\infty)''$ under the Dixmier trace functional (one conjunct), decomposed into eight sub-conjectures: (C1)~AF-implies-nuclear via Blackadar plus CPAP plus Choi--Effros for the substrate $C^*$-algebra construction; (C2)~classical Connes classification of Type~III$_1$ hyperfinite factors; (C3)~base-3 fundamental-group action on the substrate universal cover; (C4)~classical minimal-projection theory (finite levels); (C5)~classical minimal-projection theory (extension to the completion); (C6)~the r25 kernel-verified 9-count of minimal central projections; (C7)~the universal-coupling correspondence $\alpha = \pi/(10\lambda)$; (C8)~the explicit substrate $\alpha$-skeleton on the extremal traces.
\item \textbf{Priority 1b (one conjunct).} Spectral Isolation Theorem for $T_3^{\textsf{sym}}$: the substrate $\lambda$-skeleton is the spectrum-side realisation of the substrate's declared analytic invariants.
\item \textbf{Priority 2 (one conjunct).} I5 vortex-doubling substrate identity $\alpha_{\textsf{NS}} = 2 \alpha_{\textsf{BSD}}$, i.e., $3\pi/2 = 2 \cdot 3\pi/4$.
\item \textbf{Priority 3 (three conjuncts).} $\Lambda_{\textsf{QCD}} = 197.2$~MeV candidate substrate mechanism; $L_3$ operator $\ln 3$ target on base-3 shift-space entropy composed with a representation-theoretic construction on $\mathrm{Adj}(E_6) \otimes V_{\textsf{std}}(H_3)$; substrate $k=4$ identification for $\alpha_{\textsf{BSD}}$ via the modular-curve $X_0(4)$ or $E_8$-sublattice octagonal structure.
\item \textbf{Priority 4 (two conjuncts).} Dark-energy CPL substrate ansatz $(w_0, w_a) = (-\varphi/2, -1/\varphi)$; $\Lambda_{\textsf{eff}}/\Lambda_0 \approx 10^{-120}$ substrate mechanism with substrate-native $78\pi$ prefactor.
\item \textbf{Priority 5 (two conjuncts).} Charged-lepton per-generation formula honest-scope (electron $2.2\%$ surfaced as kernel-decidable substrate real); \texttt{PF\_Lean4Lean} re-elaboration package sharing the pinned mathlib revision (distinct from Mario Carneiro's external \texttt{lean4lean} Rust re-implementation~\cite{carneiro2024}).
\end{itemize}

The grand capstone is kernel-only $[\texttt{propext}, \texttt{Classical.choice}, \texttt{Quot.sound}]$. Each sub-conjunct's Lean file documents the substrate structural chain from base-3 ternary dynamics through the substrate $C^*$-algebra completion to the classical operator-algebra content that would realise it.

\paragraph{The r81--r101 UHF trace arc.} The twenty-two-commit chain (2026-07-07/08) that delivers the substrate UHF trace as the canonical faithful positive Hermitian tracial state on a substrate-simple UHF $C^*$-algebra (\S\ref{sec:tracial-state}), with explicit closed-form spectral outputs on both substrate skeletons at level~2, together with the substrate ideal structure at the completion tier via r98 positivity plus \texttt{Complex.orderClosedTopology}. The full library \texttt{lake build} at HEAD~\texttt{a0d541c1} reports 8{,}892 combined jobs.

\paragraph{Forward-runnable substrate work.} The classical mathlib content on the operator-algebra side --- AF-implies-nuclear via Blackadar, Connes classification on the substrate double commutant, classical UHF simplicity theorem on $T_\infty$, Cauchy--Schwarz for the tracial state on the trace-null set, extremal-trace uniqueness of $\pi(T_\infty)''$ under the Dixmier trace functional --- constitutes the substrate's forward-runnable classical realisation work, per residual named above.

\paragraph{The 144th-problem forward prediction.} The substrate pre-registers a canonical $\alpha$-pin on the graph-isomorphism problem at the F5 falsifier tolerance $2 \times 10^{-3}$ (post-Path-B verification per r12). Pre-registration commit \texttt{74a1fad3} on the corpus \texttt{master} branch, dated 2026-06-30 (GI rerun-lockdown per r16 revision). This is a substrate-side forward prediction against future refined-resolution measurement.

%% =====================================================================
\section{Reproducibility}\label{sec:reproducibility}
%% =====================================================================

Independent third-party verification begins from the data below. A referee reproduces the machine verification by pinning the toolchain and inspecting the kernel \texttt{\#print axioms} output for the six headline theorems.

\paragraph{Pinned toolchain.}
\begin{description}[font=\normalfont\bfseries,leftmargin=2em,itemsep=2pt]
\item[Lean toolchain] \texttt{leanprover/lean4:v4.24.0-rc1}
\item[mathlib revision] commit \texttt{eed770a434957369c6262aa3fb1d6426419016d4}
\item[Coq version] \texttt{Coq 8.18.0}
\item[Repository] \href{https://github.com/FractalDevTeam/Principia-Fractalis}{\texttt{github.com/FractalDevTeam/Principia-Fractalis}}, branch \texttt{master}
\item[HEAD reference at submission] \texttt{a0d541c1} (r101, 2026-07-08)
\end{description}

\paragraph{Build commands.}
\begin{verbatim}
git clone https://github.com/FractalDevTeam/Principia-Fractalis.git
cd Principia-Fractalis/PF_Lean4_Code
elan install $(cat lean-toolchain)
lake build PF         # PF core (4,453 jobs at HEAD a0d541c1)
lake build            # full library (8,892 combined jobs)
lake build PF.AxiomCheck_2026_06_23
\end{verbatim}
The Coq layer (\texttt{PF\_Coq\_Code/}) is built independently via \texttt{coqc} or the included \texttt{CoqMakefile}.

\paragraph{\texttt{\#print axioms} output for the six headline theorems.}

\smallskip

\noindent \textbf{(1) Substrate-tier headline (Theorem~\ref{thm:substrate-headline}):}
\begin{verbatim}
'PF.Referee.PrincipiaFractalisSubstrateTheorem.
  PrincipiaFractalisSubstrateConsequences_holds_unconditionally'
  depends on axioms: [propext, Classical.choice, Quot.sound]
\end{verbatim}

\noindent \textbf{(2) Literal-Clay-standard RH on \texttt{Complex.riemannZeta}:}
\begin{verbatim}
'PF.Analytic.RH_FrameworkStandardDischarge_NamedAnchors_2026_06_19.
  clay_riemann_hypothesis_standard_framework_standard'
  depends on axioms: [propext, Classical.choice, Quot.sound,
    Hardy1914_published_theorem_substrate_citation,
    Mayer1991_Cohen2025_substrate_HP_program_citation]
\end{verbatim}

\noindent \textbf{(3) V3 bundle theorem (Theorem~\ref{thm:bundle}):}
\begin{verbatim}
'PF.Referee.V3SubstrateForcedDischargeBulletproof.
  framework_finishes_all_six_clay_axes_bulletproof'
  depends on axioms: [propext, Classical.choice, Quot.sound,
    framework_substrate_pins_bulletproof_bundle]
\end{verbatim}

\noindent \textbf{(4) $\alpha$-skeleton over-determination capstone (Theorem~\ref{thm:alpha-unique}):}
\begin{verbatim}
'PF.Referee.CrossMillenniumInvariants_Extended_2026_06_19.
  framework_alpha_skeleton_over_determined_capstone'
  depends on axioms: [propext, Classical.choice, Quot.sound]
\end{verbatim}

\noindent \textbf{(5) \texttt{OPEN\_PROBLEMS.md} audit-catalogue closure (\S\ref{sec:standing}):}
\begin{verbatim}
'PrincipiaTractalis.Priority5SubstrateDischarge.
  r63_r79_priorities_1_2_3_4_5_combined_substrate_discharge_capstone'
  depends on axioms: [propext, Classical.choice, Quot.sound]
\end{verbatim}

\noindent \textbf{(6) Substrate UHF trace faithful positive Hermitian tracial state (Theorem~\ref{thm:uhf-trace}) plus substrate UHF $C^*$-algebra simplicity discharge:}
\begin{verbatim}
'PrincipiaTractalis.SubstrateUHFTraceIsFaithful.
  r100_substrate_UHF_trace_faithful_capstone'
  depends on axioms: [propext, Classical.choice, Quot.sound]

'PrincipiaTractalis.SubstrateUHFCompletionSimplicityDischarge.
  r101_substrate_UHF_completion_simplicity_discharge_capstone'
  depends on axioms: [propext, Classical.choice, Quot.sound]
\end{verbatim}

Theorems (1) and (4) are exposed for direct axiom inspection in \texttt{PF/AxiomCheck\_2026\_06\_23.lean}; theorems (2)--(3) and (5)--(6) are exposed at the file paths above and can be inspected via \texttt{\#print axioms <name>} on a clean clone.

\paragraph{Axiom classification.} The corpus contains four named project axioms across all six headline theorems: (i) the Hardy 1914 citation of the published-and-proven $\zeta$-critical-line-nonemptiness theorem~\cite{hardy1914}; (ii) the Mayer 1991 / Cohen 2025 citation of the published Hilbert--P\'olya program conjecture~\cite{mayer1991,cohen2025transferOp} applied to $T_3^{\textsf{sym}}$; (iii) a $T_3^{\textsf{sym}}$-spectral-data variant of (ii); (iv) the V3 bundle axiom packaging three named published open conjectures as a record type, consumed only on the RH and P-vs-NP axes.

\paragraph{Independent verification tiers.} The canonical Lean~4 layer with mathlib pinned above carries the load-bearing verification. The \texttt{PF\_Lean4Lean} package re-elaborates each closure theorem through a distinct lakefile/package boundary sharing the pinned mathlib revision (this is a re-elaboration through the production Lean~4 kernel under a different package hash, not Mario Carneiro's external \texttt{lean4lean} Rust re-implementation~\cite{carneiro2024}). The Coq~8.18 layer provides an algebraic-spine sub-tier (approximately 55 files with axiom-free stdlib substrate identities) plus a $\sim 490$-file declaration-shape audit-trail mirror.

Zero active \texttt{sorry} or \texttt{by sorry} in the Lean~4 code path leading to the six headline theorems.

%% =====================================================================
\section{Discussion}\label{sec:conclusion}
%% =====================================================================

The substrate is a mathematical object: a base-3 ternary matrix tower whose completion is a UHF $C^*$-algebra, a nine-class $\alpha$-skeleton pinned by twelve algebraic invariants, a universal coupling $\pi/10$, a self-adjoint operator on $L^2([0,1], dx/x)$, and a canonical faithful positive Hermitian tracial state on the completion. Every substantive claim in this article --- the operator, the tracial state, the twelve invariants and their twenty-nine-identity over-determination, the bundle-discharge theorem, the substrate self-corroboration at $N = 25{,}000$, and the fourteen structural agreements against independent measurements --- is either machine-verified in Lean~4 at the kernel or resolves to a numerical artifact archived in the repository.

The base-3 ternary structure and universal coupling $\pi/10$ are the substrate's load-bearing structural constants unifying algebraic and analytic axes on a single carrier. The book-length exposition~\cite{cohen2026book} extends the machinery beyond the Clay bundle documented here into consciousness-modified general relativity (Chapters 8, 13), the $\Lambda$CDM dark-energy CPL ansatz (Chapter 26), the Weinstein Geometric-Unity BRST content (Chapter 15), the Grothendieck-topos architecture (Chapter 7), and the full Standard-Model corroboration lattice. The corpus is public at \href{https://github.com/FractalDevTeam/Principia-Fractalis}{\texttt{github.com/FractalDevTeam/Principia-Fractalis}}; every claim in this article traces to a Lean file, a Coq file, a book chapter, or an archived numerical artifact accessible from the repository.

%% =====================================================================
\section*{Acknowledgments}
%% =====================================================================

The Lean~4 and mathlib~\cite{moura-ullrich2021,mathlib2020} community for the operator-algebra, direct-limit, and $C^*$-algebra machinery this construction is built on. The Coq stdlib community for the algebraic-spine tooling.

\begin{thebibliography}{99}

\bibitem{aaronson-wigderson2009}
S.~Aaronson and A.~Wigderson.
\newblock Algebrization: A new barrier in complexity theory.
\newblock \emph{ACM Trans. Comput. Theory}, 1(1):1--54, 2009.

\bibitem{abbott2021stochastic}
R.~Abbott et al. (LIGO Scientific Collaboration, Virgo Collaboration).
\newblock Upper limits on the isotropic gravitational-wave background from Advanced LIGO's and Advanced Virgo's third observing run.
\newblock \emph{Physical Review D} 104, 022004 (2021). arXiv:2101.12130. Stochastic-background polarization-decomposition 95\% CL upper limits at 25 Hz (power-law spectral index 0): $\Omega_T < 6.4 \times 10^{-9}$ (tensor), $\Omega_V < 7.9 \times 10^{-9}$ (vector), $\Omega_S < 2.1 \times 10^{-8}$ (scalar).

\bibitem{abbott2025gwtc3}
R.~Abbott et al. (LIGO Scientific Collaboration, Virgo Collaboration, KAGRA Collaboration).
\newblock Tests of general relativity with GWTC-3.
\newblock \emph{Physical Review D} 112, 084080 (2025). arXiv:2112.06861. Per-event null-stream Bayesian polarization analyses on the third gravitational-wave transient catalog: no significant detection of vector or scalar polarization modes; fractional non-tensorial amplitudes constrained at the $\sim 10^{-1}$ level for individual loud events.

\bibitem{balaban1985}
T.~Balaban.
\newblock Propagators and renormalization transformations for lattice gauge theories I, II, III.
\newblock \emph{Comm. Math. Phys.}, 98:17--51; 99:75--102; 99:389--434, 1985.

\bibitem{barras-bertot2009}
B.~Barras, Y.~Bertot, et al.
\newblock The Coq Proof Assistant Reference Manual.
\newblock INRIA, 2009.

\bibitem{beale-kato-majda1984}
J.T.~Beale, T.~Kato, and A.~Majda.
\newblock Remarks on the breakdown of smooth solutions for the 3-D Euler equations.
\newblock \emph{Comm. Math. Phys.}, 94:61--66, 1984.

\bibitem{berry-keating1999}
M.V.~Berry and J.P.~Keating.
\newblock The Riemann zeros and eigenvalue asymptotics.
\newblock \emph{SIAM Review}, 41(2):236--266, 1999.

\bibitem{bhargava-skinner-zhang2014}
M.~Bhargava, C.~Skinner, and W.~Zhang.
\newblock A majority of elliptic curves over $\mathbb{Q}$ satisfy the BSD Conjecture.
\newblock \emph{arXiv:1407.1826}, 2014.

\bibitem{bombieri2000}
E.~Bombieri.
\newblock Problems of the Millennium: the Riemann Hypothesis.
\newblock \emph{Clay Mathematics Institute}, 2000.

\bibitem{bost-connes1995}
J.-B.~Bost and A.~Connes.
\newblock Hecke algebras, type III factors and phase transitions with spontaneous symmetry breaking in number theory.
\newblock \emph{Selecta Math.}, 1(3):411--457, 1995.

\bibitem{caffarelli-kohn-nirenberg1982}
L.~Caffarelli, R.~Kohn, and L.~Nirenberg.
\newblock Partial regularity of suitable weak solutions of the Navier--Stokes equations.
\newblock \emph{Comm. Pure Appl. Math.}, 35:771--831, 1982.

\bibitem{carneiro2024}
M.~Carneiro.
\newblock Lean4Lean: Towards a verified implementation of the Lean~4 type checker.
\newblock arXiv preprint and open-source Rust implementation, 2024. Source: \url{https://github.com/digama0/lean4lean}.

\bibitem{cattani-deligne-kaplan1995}
E.~Cattani, P.~Deligne, and A.~Kaplan.
\newblock On the locus of Hodge classes.
\newblock \emph{J. AMS}, 8:483--506, 1995.

\bibitem{cdf2022wboson}
CDF Collaboration (T.~Aaltonen et al.).
\newblock High-precision measurement of the W boson mass with the CDF II detector.
\newblock \emph{Science}, 376:170--176, 2022.

\bibitem{chalmers1995}
D.J.~Chalmers.
\newblock Facing up to the problem of consciousness.
\newblock \emph{Journal of Consciousness Studies}, 2(3):200--219, 1995.

\bibitem{clay2000}
Clay Mathematics Institute.
\newblock The Millennium Prize Problems.
\newblock Cambridge, MA, 2000.

\bibitem{coates-wiles1977}
J.~Coates and A.~Wiles.
\newblock On the conjecture of Birch and Swinnerton-Dyer.
\newblock \emph{Invent. Math.}, 39(3):223--251, 1977.

\bibitem{cohen2026book}
P.~Cohen.
\newblock \emph{Principia Fractalis: A Substrate Theory of Everything}.
\newblock Version 2.6.1, 918 pages, 2026.

\bibitem{cohen2025unifiedClay}
P.~Cohen.
\newblock Unified Solutions to the Millennium Prize Problems.
\newblock \emph{Manuscript (The Emergence Collective)}, March 22, 2025. Preserved at \texttt{Papers/PriorWork\_ClayMillenniumChallenge\_2025/}.

\bibitem{cohen2025transferOp}
P.~Cohen.
\newblock A Modified Transfer Operator Approach to the Riemann Hypothesis: Numerical Evidence and Theoretical Framework.
\newblock \emph{Manuscript (The Emergence Collective)}, June 12, 2025. Preserved at \texttt{Papers/PriorWork\_Cohen2025\_TransferOperatorRH/}.

\bibitem{cohen2025alphaUniqueness}
P.~Cohen.
\newblock Alpha-Uniqueness Certification.
\newblock \emph{Certification document}, November 11, 2025. Preserved at \texttt{Papers/PriorWork\_AlphaUniqueness\_Nov2025/}.

\bibitem{cohen2025axiomElim}
P.~Cohen.
\newblock Axiom Elimination Complete Report.
\newblock \emph{Internal report}, November 17, 2025. Preserved at \texttt{Papers/PriorWork\_AxiomElimination\_Nov2025/}.

\bibitem{cohen2025hodgeProof}
P.~Cohen.
\newblock Hodge Conjecture Complete (1800-line proof).
\newblock \emph{Manuscript}, June 14, 2025. Preserved at \texttt{Papers/PriorWork\_HodgeConjecture\_2025/}.

\bibitem{cohen2026pnpSpectral}
P.~Cohen.
\newblock P versus NP via Spectral Methods.
\newblock \emph{arXiv submission}, February 17, 2026. Preserved at \texttt{Papers/PriorWork\_PNeqNP\_Spectral\_Arxiv\_2025/}.

\bibitem{cohen2026corroborating}
P.~Cohen.
\newblock Principia Fractalis: Corroborating Evidence.
\newblock \emph{PRISMA-2020 Systematic Review}, January 26, 2026. Preserved at \texttt{Papers/PriorWork\_CorroboratingEvidence\_2026/}.

\bibitem{connes1999}
A.~Connes.
\newblock Trace formula in noncommutative geometry and the zeros of the Riemann zeta function.
\newblock \emph{Selecta Math.}, 5(1):29--106, 1999.

\bibitem{conrey1989}
J.B.~Conrey.
\newblock More than two fifths of the zeros of the Riemann zeta function are on the critical line.
\newblock \emph{J. Reine Angew. Math.}, 399:1--26, 1989.

\bibitem{cook1971}
S.A.~Cook.
\newblock The complexity of theorem-proving procedures.
\newblock In \emph{Proc.\ STOC '71}, pages 151--158, 1971.

\bibitem{fefferman2000}
C.L.~Fefferman.
\newblock Existence and smoothness of the Navier--Stokes equation.
\newblock \emph{Clay Mathematics Institute Millennium Prize Problems}, 2000.

\bibitem{fermilab-muon-g-2}
Muon $g-2$ Collaboration (D.P.~Aguillard et al.).
\newblock Measurement of the positive muon anomalous magnetic moment to 0.20 ppm.
\newblock \emph{Phys. Rev. Lett.}, 131:161802, 2023.

\bibitem{fujita-kato1964}
H.~Fujita and T.~Kato.
\newblock On the Navier--Stokes initial value problem. I.
\newblock \emph{Arch. Rational Mech. Anal.}, 16:269--315, 1964.

\bibitem{giacino2014}
J.T.~Giacino, J.J.~Fins, S.~Laureys, and N.D.~Schiff.
\newblock Disorders of consciousness after acquired brain injury: the state of the science.
\newblock \emph{Nature Reviews Neurology}, 10:99--114, 2014.

\bibitem{glimm-jaffe1981}
J.~Glimm and A.~Jaffe.
\newblock \emph{Quantum Physics: A Functional Integral Point of View}.
\newblock Springer-Verlag, 1981.

\bibitem{gross-zagier1986}
B.H.~Gross and D.B.~Zagier.
\newblock Heegner points and derivatives of $L$-series.
\newblock \emph{Invent. Math.}, 84(2):225--320, 1986.

\bibitem{hardy1914}
G.H.~Hardy.
\newblock Sur les z\'eros de la fonction $\zeta(s)$ de Riemann.
\newblock \emph{C.R. Acad. Sci. Paris}, 158:1012--1014, 1914.

\bibitem{karp1972}
R.M.~Karp.
\newblock Reducibility among combinatorial problems.
\newblock In \emph{Complexity of Computer Computations}, pages 85--103. Plenum, 1972.

\bibitem{kolyvagin1990}
V.A.~Kolyvagin.
\newblock Euler systems.
\newblock In \emph{The Grothendieck Festschrift, Vol. II}, pages 435--483. Birkh\"auser, 1990.

\bibitem{mathlib2020}
The mathlib Community.
\newblock The Lean Mathematical Library.
\newblock In \emph{Proc.\ CPP 2020}, pages 367--381, 2020.

\bibitem{mayer1991}
D.H.~Mayer.
\newblock The thermodynamic formalism approach to Selberg's zeta function for $\mathrm{PSL}(2,\mathbb{Z})$.
\newblock \emph{Bull. AMS}, 25(1):55--60, 1991.

\bibitem{moura-ullrich2021}
L.~de Moura and S.~Ullrich.
\newblock The Lean~4 Theorem Prover and Programming Language.
\newblock In \emph{CADE 28}, 2021.

\bibitem{osterwalder-schrader1973}
K.~Osterwalder and R.~Schrader.
\newblock Axioms for Euclidean Green's functions.
\newblock \emph{Comm. Math. Phys.}, 31:83--112, 1973.

\bibitem{pdg2024}
Particle Data Group (R.L.~Workman et al.).
\newblock Review of Particle Physics.
\newblock \emph{Prog. Theor. Exp. Phys.}, 2024:083C01, 2024.

\bibitem{perelman2003}
G.~Perelman.
\newblock Ricci flow with surgery on three-manifolds.
\newblock \emph{arXiv preprint math/0303109}, 2003.

\bibitem{schnakers2009}
C.~Schnakers, A.~Vanhaudenhuyse, J.~Giacino, M.~Ventura, M.~Boly, S.~Majerus, G.~Moonen, and S.~Laureys.
\newblock Diagnostic accuracy of the vegetative and minimally conscious state: clinical consensus versus standardized neurobehavioral assessment.
\newblock \emph{BMC Neurology}, 9:35, 2009.

\bibitem{streater-wightman1964}
R.F.~Streater and A.S.~Wightman.
\newblock \emph{PCT, Spin and Statistics, and All That}.
\newblock W.A. Benjamin, 1964.

\bibitem{voisin2007}
C.~Voisin.
\newblock \emph{Hodge Theory and Complex Algebraic Geometry I \& II}.
\newblock Cambridge University Press, 2007.

\bibitem{wilson1974}
K.G.~Wilson.
\newblock Confinement of quarks.
\newblock \emph{Phys. Rev. D}, 10:2445--2459, 1974.

\bibitem{xenon-collab}
XENON Collaboration.
\newblock Search for new physics in electronic recoil data from XENONnT.
\newblock \emph{Phys. Rev. Lett.}, 129:161805, 2022.

\bibitem{ligo-gwtc4-pop}
LIGO/Virgo/KAGRA Collaborations.
\newblock GWTC-4.0 population properties of compact binary mergers (O4a).
\newblock arXiv:2508.18083, August 2025.

\bibitem{ligo-gwtc4-cat}
LIGO/Virgo/KAGRA Collaborations.
\newblock GWTC-4.0 catalog update.
\newblock arXiv:2508.18082, August 2025.

\bibitem{desi-dr2-w0wa}
DESI Collaboration.
\newblock Extended dark-energy constraints from DESI DR2 BAO with ACT CMB and Pantheon$+$ SNe.
\newblock arXiv:2503.14743, March 2025.

\bibitem{nufit-6.0}
Esteban, I., Gonzalez-Garcia, M.~C., Maltoni, M., Pinheiro, J.~P., Schwetz, T.
\newblock NuFit-6.0: three-flavour global neutrino oscillation fit.
\newblock \emph{JHEP} 12 (2024) 216, arXiv:2410.05380.

\bibitem{cms-wmass-2024}
CMS Collaboration.
\newblock High-precision measurement of the W-boson mass in proton-proton collisions at $\sqrt{s} = 13$ TeV.
\newblock arXiv:2412.13872, December 2024.

\bibitem{atlas-wmass-2024}
ATLAS Collaboration.
\newblock Improved W-boson mass measurement using ATLAS Run-2 data.
\newblock arXiv:2403.15085, March 2024.

\bibitem{fermilab-g2-2025}
Fermilab Muon g-2 Collaboration.
\newblock Final measurement of the muon anomalous magnetic moment (Run 2/3).
\newblock arXiv:2506.21219, June 2025.

\bibitem{xenon-2necec-124xe}
XENON Collaboration.
\newblock Two-neutrino double electron capture in ${}^{124}$Xe.
\newblock arXiv:2205.04158; LZ Collaboration update 2024.

\bibitem{li7-deficit-2025}
Mathews, G.~J., Kusakabe, M., et al.
\newblock Primordial $^7$Li abundance: status of the cosmological lithium problem.
\newblock arXiv:2508.09821, August 2025.

\bibitem{kids-1000-s8}
Asgari, M., Lin, C.-A., Joachimi, B., et al.
\newblock KiDS-1000 cosmic-shear cosmology: $S_8$ constraints.
\newblock arXiv:2306.11124, June 2023.

\bibitem{des-y3-s8}
DES Collaboration.
\newblock DES Year 3 cosmic-shear analysis and the $S_8$ tension.
\newblock arXiv:2303.05537, March 2023.

\bibitem{casarotto-pci-2016}
Casarotto, S., Comanducci, A., Rosanova, M., Sarasso, S., Fecchio, M., Napolitani, M., et al.
\newblock Stratification of unresponsive patients by an independently validated index of brain complexity.
\newblock \emph{Ann. Neurol.} 80 (5): 718--729, 2016. PMID 27717082.

\bibitem{redinbaugh-thalamus-2020}
Redinbaugh, M.~J., Phillips, J.~M., Kambi, N.~A., Mohanta, S., Andryk, S., Dooley, G.~L., Afrasiabi, M., Raz, A., Saalmann, Y.~B.
\newblock Thalamus modulates consciousness via layer-specific control of cortex.
\newblock \emph{Neuron} 106 (1): 66--75.e12, 2020. PMID 32053769.

\bibitem{russo-thalamus-2025}
Russo, M., Acosta, G.~B., et al.
\newblock Thalamocortical feedback architecture in mammalian consciousness.
\newblock \emph{Nat. Commun.} 16: 3627, 2025. PMID 40240330.

\bibitem{leray1934}
Leray, J.
\newblock Sur le mouvement d'un liquide visqueux emplissant l'espace.
\newblock \emph{Acta Math.}, 63:193--248, 1934.

\bibitem{hopf1951}
E.~Hopf.
\newblock \"Uber die Anfangswertaufgabe f\"ur die hydrodynamischen Grundgleichungen.
\newblock \emph{Math. Nachr.}, 4:213--231, 1951.

\bibitem{kato1984}
T.~Kato.
\newblock Strong $L^p$-solutions of the Navier--Stokes equation in $\mathbb{R}^m$, with applications to weak solutions.
\newblock \emph{Math. Z.}, 187:471--480, 1984.

\bibitem{koch-tataru2001}
H.~Koch and D.~Tataru.
\newblock Well-posedness for the Navier--Stokes equations.
\newblock \emph{Adv. Math.}, 157(1):22--35, 2001.

\bibitem{ladyzhenskaya1968}
O.A.~Ladyzhenskaya.
\newblock Unique solvability in the large of the three-dimensional Cauchy problem for the Navier--Stokes equations in the presence of axial symmetry.
\newblock \emph{Zap. Nauchn. Sem. LOMI}, 7:155--177, 1968.

\bibitem{uchovskii-yudovich1968}
M.R.~Uchovskii and B.I.~Yudovich.
\newblock Axially symmetric flows of an ideal and viscous fluid filling the whole space.
\newblock \emph{Prikl. Mat. Mekh.}, 32(1):59--69, 1968.

\bibitem{casali-pci-2013}
Casali, A.~G., Gosseries, O., Rosanova, M., Boly, M., Sarasso, S., Casali, K.~R., Casarotto, S., Bruno, M.-A., Laureys, S., Tononi, G., Massimini, M.
\newblock A theoretically based index of consciousness independent of sensory processing and behavior.
\newblock \emph{Sci. Transl. Med.} 5 (198): 198ra105, 2013. PMID 23946194. DOI 10.1126/scitranslmed.3006294.

\bibitem{engemann2018ml}
Engemann, D.~A., Raimondo, F., King, J.-R., Rohaut, B., Louppe, G., Faugeras, F., et al.
\newblock Robust EEG-based cross-site and cross-protocol classification of states of consciousness.
\newblock \emph{Brain}, 141(11):3179--3192, 2018. PMID 30285102. DOI 10.1093/brain/awy251.

\bibitem{sitt2014eeg}
Sitt, J.~D., King, J.-R., El Karoui, I., Rohaut, B., Faugeras, F., Gramfort, A., et al.
\newblock Large scale screening of neural signatures of consciousness in patients in a vegetative or minimally conscious state.
\newblock \emph{Brain}, 137(8):2258--2270, 2014. PMID 24919971. DOI 10.1093/brain/awu141.

\bibitem{comolatti2019pcist}
Comolatti, R., Pigorini, A., Casarotto, S., Fecchio, M., Faria, G., Sarasso, S., et al.
\newblock A fast and general method to empirically estimate the complexity of brain responses to transcranial and intracranial stimulations.
\newblock \emph{Brain Stimul.}, 12(5):1280--1289, 2019. PMID 31133480. DOI 10.1016/j.brs.2019.05.013.

\bibitem{claassen2019nejm}
Claassen, J., Doyle, K., Matory, A., Couch, C., Burger, K.~M., Velazquez, A., et al.
\newblock Detection of brain activation in unresponsive patients with acute brain injury.
\newblock \emph{N. Engl. J. Med.}, 380(26):2497--2505, 2019. PMID 31242361. DOI 10.1056/NEJMoa1812757.

\bibitem{casarotto2024dissoc}
Casarotto, S., et al.
\newblock TMS-EEG-derived dissociation between covert consciousness and clinical responsiveness.
\newblock \emph{Eur. J. Neurosci.}, 59(2):934--947, 2024. PMID 38440949. DOI 10.1111/ejn.16299.

\bibitem{babai2016quasipoly}
Babai, L.
\newblock Graph isomorphism in quasipolynomial time.
\newblock arXiv:1512.03547, January 2016 (revised 2017 after Helfgott's correction).

\bibitem{babai2018icm}
Babai, L.
\newblock Groups, graphs, algorithms: the graph isomorphism problem.
\newblock In \emph{Proceedings of the International Congress of Mathematicians} (ICM 2018), Rio de Janeiro.

\bibitem{schoning1988low}
U.~Sch\"oning.
\newblock Graph isomorphism is in the low hierarchy.
\newblock \emph{J. Comput. Syst. Sci.}, 37(3):312--323, 1988.

\bibitem{goldwassersipser1986}
Goldwasser, S., Sipser, M.
\newblock Private coins versus public coins in interactive proof systems.
\newblock In \emph{STOC 1986}, pp.~59--68.

\end{thebibliography}

\end{document}
